\documentclass[11pt]{article}
\usepackage[a4paper,margin=28mm]{geometry}
\usepackage[T1]{fontenc}
\usepackage{lmodern,microtype,amsmath,amssymb,amsthm}
\usepackage[hidelinks]{hyperref}
\newtheorem{theorem}{Theorem}
\newtheorem{lemma}[theorem]{Lemma}
\newcommand{\tr}{\operatorname{tr}}
\newcommand{\diag}{\operatorname{diag}}
\title{The one-column spectral minimax identity\\[3pt]\large A special case of TR-03, not a full resolution}
\author{Research note prepared for mathematical review}
\date{12 September 2026}
\begin{document}
\maketitle
\noindent\textbf{Status.} The argument below proves the $k=1$ case for every positive spectrum. It does not determine the joint dependence on $(n,k)$ requested in TR-03. This elementary identity may already be implicit in earlier work; no priority or novelty claim is made.

\section{Statement}
Let $n\ge3$ and $\lambda=(\lambda_1,\ldots,\lambda_n)\in(0,\infty)^n$. Put $\Lambda=\diag(\lambda)$ and $K_V=V^T\Lambda V$ for $V\in O(n)$. Define
\[
 x_1(\lambda)=\max_{V\in O(n)}\min_{1\le i\le n}
 \tr\left(K_V-\frac{(K_V)_{:i}(K_V)_{i:}}{(K_V)_{ii}}\right),
 \qquad y_1(\lambda)=\frac{2e_2(\lambda)}{e_1(\lambda)}.
\]
These are the $k=1$ quantities in the repository target~\cite{repo}.
\begin{theorem}
For every positive spectrum,
\[
 x_1(\lambda)=y_1(\lambda)
 =\sum_{j=1}^n\lambda_j-
   \frac{\sum_{j=1}^n\lambda_j^2}{\sum_{j=1}^n\lambda_j}.
\]
Consequently $R_{n,1}=\sup_{\lambda>0}y_1(\lambda)/x_1(\lambda)=1$.
\end{theorem}

\section{A zero-diagonal orthogonal basis}
\begin{lemma}
Every real symmetric matrix $H$ of trace zero has an orthonormal basis $v_1,\ldots,v_n$ satisfying $v_i^THv_i=0$ for every $i$.
\end{lemma}
\begin{proof}
Induct on the dimension. The one-dimensional case is immediate. If $H=0$, every basis works. Otherwise trace zero implies the presence of a positive and a negative eigenvalue. A unit vector $v_1$ in the span of corresponding eigenvectors has zero Rayleigh quotient by continuity (or by choosing the squared component weights inversely proportional to the absolute eigenvalues). The compression of $H$ to $v_1^\perp$ is symmetric and has trace $\tr H-v_1^THv_1=0$. Apply induction to that compression and append $v_1$.
\end{proof}

\section{Proof of the identity}
Write $S_1=\sum_j\lambda_j$, $S_2=\sum_j\lambda_j^2$ and $c=S_2/S_1$. For any orthogonal $V$ and $K=K_V$, the error of choosing column $i$ is
\[
 E_i(K)=S_1-\frac{(K^2)_{ii}}{K_{ii}}.
\]
The weights $K_{ii}/S_1$ are positive and sum to one. Hence
\[
 \min_i E_i(K)\le
 \sum_i\frac{K_{ii}}{S_1}E_i(K)
 =S_1-\frac{S_2}{S_1}=y_1(\lambda).
\]
This proves $x_1\le y_1$.

For the reverse inequality, set $H=\Lambda^2-c\Lambda$. Its trace is zero. By the lemma, choose an orthogonal matrix $V$ whose columns have zero $H$-Rayleigh quotient. Since $(V^T\Lambda V)^2=V^T\Lambda^2V$, the resulting $K$ obeys
\[
 (K^2)_{ii}=cK_{ii}\qquad(1\le i\le n).
\]
Thus \emph{every} column has the same error $E_i(K)=S_1-c=y_1(\lambda)$. Therefore $x_1\ge y_1$, proving equality. The formula for $y_1$ follows from $2e_2=S_1^2-S_2$. Positivity of the spectrum and $n\ge3$ ensure $y_1>0$, so the ratio is well defined.

\section{What remains}
For $2\le k\le n-2$, neither this argument nor the accompanying numerical examples determine the sharp spectral minimax ratio. Simultaneous equalization of all $k$-column residuals is substantially different from equalizing the diagonal of a single trace-zero symmetric matrix. This note must not be submitted or counted as a resolution of the full TR-03 target.

The script \texttt{verification/verify\_TR03.py} checks one radical-valued example exactly and constructs zero-diagonal bases for a fixed set of numerical spectra. These computations are illustrative; the proof above is independent of them.

\begin{thebibliography}{9}
\bibitem{repo} OpenProblemsInNLA, \emph{TR-03: Sharp gap between volume sampling and the worst matrix with a prescribed spectrum}, accessed 12 September 2026.\\
\url{https://github.com/ajt60gaibb/OpenProblemsInNLA/tree/main/randomized-and-low-rank-approximation/TR-03}
\end{thebibliography}
\end{document}
